\section{The maximum diagonal-stable smooth class}

\paragraph{I.1 (Proved).}
Let $p:E\to B$ be a functor of ordinary small categories, and form
its strict diagonal $\Delta_p:E\to E\times_B E$. Then
\[
 p\text{ and }\Delta_p\text{ are }W_\infty\text{-smooth}
 \quad\Longleftrightarrow\quad
 p\text{ is a discrete fibration}.
\]
Consequently the discrete fibrations form the largest subclass of
$\operatorname{Liss}_{W_\infty}$ closed under diagonals. They are also
closed under identities, composition and strict pullback.
\begin{proof}
For smoothness use the incoming lifting category: given $x\in E$
and $a:b\to p(x)$, its objects are arrows $h:y\to x$ with
$p(y)=b$ and $p(h)=a$, and its arrows are the commuting triangles
over $1_b$. If $p$ is smooth, this category is weakly contractible,
and in particular nonempty.

Take any two such lifts $h:y\to x$ and $k:z\to x$. They give an
arrow $(h,k):(y,z)\to(x,x)$ in the strict pullback. A lift of this
arrow through $\Delta_p$ would be an arrow $l:w\to x$ satisfying
$(w,w)=(y,z)$ and $(l,l)=(h,k)$. Thus its incoming lifting category
is empty unless $y=z$ and $h=k$. Smoothness of $\Delta_p$ excludes
emptiness. Existence and uniqueness of each incoming lift is exactly
the discrete-fibration condition.

Conversely suppose $p$ is a discrete fibration. Its incoming lifting
categories are singletons. For an arrow $(h,k)$ into a diagonal
object, the discrete-fibration uniqueness gives $h=k$ and $y=z$;
the diagonal lifting category is again a singleton. Therefore both
maps are smooth. The usual unique-lift argument proves stability
under composition and strict pullback. Any diagonal-stable subclass
of the smooth maps satisfies the left side of the equivalence for
each of its maps and is therefore contained in this class.
\end{proof}


\paragraph{I.2 (Proved).}
In the presheaf Kan formalism with space coefficients, a discrete
fibration $p:E\to B$ has universally equivalent left and right Kan
extensions if and only if $p$ is an isomorphism. In particular the
stable classification by finite discrete bifibrations does not carry
over to unpointed spaces.
\begin{proof}
Pull back to any object $b\in B$. The fibre is a discrete set $I$.
Left and right Kan extension to a point are respectively the
coproduct and product functors $\mathcal S^I\to\mathcal S$.
On the family constantly equal to $*$ they yield the discrete set
$I$ and $*$, respectively. Hence equivalence forces $I$ to contain
exactly one element. In a discrete fibration with exactly one
object over each object of $B$, every arrow of $B$ has exactly one
lift with its specified target. This makes $p$ bijective on objects
and on all hom sets, and therefore an isomorphism. Conversely an
isomorphism has equal inverse, left-Kan and right-Kan functors,
universally under strict pullback.
\end{proof}


\paragraph{I.3 (Proved).}
Let $\mathcal V$ be a presentable closed symmetric monoidal
infinity-category with tensor product preserving small colimits in
each variable. The assignment
$\Phi(S)=\operatorname{Fun}(NS^{\mathrm{op}},\mathcal V)$ extends to
a lax symmetric monoidal correspondence functor with smooth right
legs, sending $X\xleftarrow aM\xrightarrow bY$ to $b_!a^*$.
This independently completes the coherence argument in BISITE F.10.
\begin{proof}
We use the precise universal property of spans in \citet[Corollary 2.6(2)]{GHIJCnossenLenzRamzi2025}: for an infinity-category $C$
and a pullback-stable wide subcategory $E$, restriction from the
span $(\infty,2)$-category with arbitrary left legs and right legs
in $E$ is an equivalence onto contravariant functors on $C$ whose
restrictions along $E$ have left adjoints and satisfy Beck--Chevalley.
The equivalence includes adjointable transformations and their
higher coherences. No diagonal or cancellation hypothesis occurs
in this statement.

Take $C=N\operatorname{Cat}$, so its pullbacks are the ordinary
strict categorical pullbacks. The incoming-lift characterization
of smoothness shows directly that smooth maps form a wide
subcategory closed under pullback and composition. Restriction of
presheaves is a coherent functor on $C^{\mathrm{op}}$, and all left
Kan adjoints exist by presentability. If $b:M\to Y$ is smooth,
the inclusion
$M_y^{\mathrm{op}}\to(y\backslash b)^{\mathrm{op}}$ is final:
its relevant comma categories are exactly the incoming lifting
categories. Hence
\[
 (b_!A)(y)\simeq\operatorname*{colim}_{M_y^{\mathrm{op}}}A.
\]
In every strict pullback square the fibre categories on the two
sides are identical. Applying this formula identifies the canonical
Beck--Chevalley mate as an equivalence. All hypotheses of the
quoted universal property now hold. It produces the correspondence
functor, and its value on a span is the stated composite.

For completeness the monoidal enhancement is also covered by
this universal property, and does not require choosing separate
associativity isomorphisms. For each $k$, external tensor is the
natural transformation
\[
 \prod_{i=1}^k\Phi(S_i)\longrightarrow
 \Phi\!\left(\prod_{i=1}^kS_i\right),\qquad
 (A_i)\longmapsto\boxtimes_i A_i.
\]
It is left-adjointable in every smooth variable. Indeed Fubini
for colimits and separate cocontinuity identify, with their
canonical units,
\[
 (\textstyle\prod_i b_i)_!(\boxtimes_i A_i)
       \simeq\boxtimes_i (b_i)_!A_i.
\]
Apply the transformation part of the span universal property to
$C^k$, with right class $E^k$. Products commute with strict
pullbacks, so both sides define functors on the corresponding
span category. The displayed transformation has a unique coherent
extension. The original associativity, symmetry and unit
coherences of external tensor lift through the fully faithful
restriction equivalence, including all higher relations. These
extensions are exactly a lax symmetric monoidal structure.

Right Kan adjoints to restrictions exist, $b_!$ has right adjoint
$b^*$, and pointwise tensor is closed by presentability and
separate cocontinuity. These are the remaining adjoints in the
2022 definition of a six-functor formalism. The projection formula
can also be checked directly by the fibre formula: restriction
of $b^*M$ to $M_y$ is constant at $M(y)$, and tensoring with this
object commutes with the displayed colimit.
\end{proof}


\paragraph{I.4 (Proved).}
The preceding proof does not follow from \citet[Proposition A.5.10]{GHIJMann2022} by taking $I$ to be all smooth maps and $P$ to be the
isomorphisms. That specialization violates the suitable-decomposition
cancellation requirement. Indeed $p:[1]\to*$ and $pg:*\to*$,
where $g$ selects $1\in[1]$, are smooth, but $g$ is not smooth:
the incoming arrow $0\to1$ has no lift. Thus $p,pg\in I$ does
not imply $g\in I$. The span universal property used in I.3 is
the needed replacement. Contractible uniqueness of individual
adjoints, without this extension statement or an explicit
unfurling construction, by itself does not state the necessary
coherence theorem. The program
\texttt{code/GHIJ/small\_counterexamples.py} enumerates the
incoming lifts for $p$, $pg$, and $g$, verifying this finite
cancellation counterexample.


\paragraph{I.5 (Proved).}
Opposition of the base categories transports I.3 to a correspondence
formalism with proper right legs and values
$\Psi(S)=\operatorname{Fun}(NS,\mathcal V)$. The maximum
diagonal-stable proper class is the discrete opfibrations. All
allowed maps of either Kan formalism satisfy $f^!=f^*$, so they
are cohomologically etale with trivial twist. These assertions
follow by applying $S\mapsto S^{\mathrm{op}}$ to I.1 and I.3;
incoming lifting categories become opposites of outgoing lifting
categories, and opposition preserves weak contractibility.


\paragraph{I.6 (Proved).}
For pointed-space coefficients, the analogous universally
ambidextrous discrete fibrations are exactly the inclusions of
unions of connected components, up to isomorphism over the base.
\begin{proof}
At a point fibre $I$, left and right Kan extension are wedge and
product. Applied to $S^0$ in every coordinate, their numbers of
components are $1+|I|$ and $2^{|I|}$. They differ when $|I|\geq2$,
including infinite cardinalities by Cantor's theorem. Thus each
fibre is empty or a singleton. Pull back along a base arrow
$0\to1$. A transition from a nonempty fibre over $0$ to an
empty one over $1$ gives the discrete fibration $\{0\}\to[1]$.
On a nonzero pointed coefficient its left Kan value at $1$ is
zero and its right Kan value is that coefficient. The reverse
emptiness pattern is excluded already by the existence of the
discrete-fibration reindexing function. Consequently the support
is a union of components, and all transitions there are bijections.
Conversely extension by zero on such a union is both Kan
extensions, universally after pullback.
\end{proof}


\paragraph{I.7 (Open).}
No proof or counterexample is supplied here for the reversible
accessible-localizer classification in Part I(2), including its
interval between integral homology equivalences and $W_0$.
For Part I(3), the unresolved assertion is an assignment
$\mathbb D\mapsto\mathbb D'$ on all derivators satisfying
Der~1--4, natural in exact morphisms, with
$W_{\mathbb D'}=W_{\mathbb D}^{\mathrm{lc}}$; we also do not
supply the alternative obstruction by two derivators equivalent
on finite posets but with different locally constant towers.
The weakest sufficient enhancement and an iteration theorem in
that generality remain unresolved. I.9 below proves a specified
localized-space realization and its finite iteration for the
enhanced categories of spectra and $D(\mathbb F_p)$, including
both second-stage calculations. It does not resolve the
bare-derivator assertion or the reversible-localizer classification.


\paragraph{I.8 (Proved).}
For spectra or $D(k)$, with $k$ a nonzero field, the universally
ambidextrous discrete fibrations are precisely the finite
discrete bifibrations: their fibres are finite sets, and all
reindexing functions along base arrows are bijections. Together
with I.2 and I.6 this gives the unstable and stable proper-map
classifications within the maximum diagonal-stable class.
Every map in that class is cohomologically smooth and etale,
with trivial dualizing twist, because $f_!$ is left adjoint to
$f^*$ and hence $f^!=f^*$.
\begin{proof}
At a point fibre $I$, ambidexterity equates coproduct and product.
An infinite product cannot preserve all coproducts: take the
family of diagrams with the coefficient unit in one coordinate
and zero in every other coordinate. The comparison is the
canonical map from the direct sum to the product of the
nonzero groups $\mathbb Z$ or $k$, and it is not surjective
for an infinite index set. Compactness of the sphere or of
$k$ detects this failure. Thus every fibre is finite.

Pull back along an arrow $0\to1$, and let
$r:I_1\to I_0$ be its reindexing map. If $i\in I_0$ has no
preimage, put a nonzero coefficient $M$ at $i$ and zero
elsewhere. Its left Kan value at $1$ is zero, while its right
Kan value contains $M$ as the sole nonzero factor. If $i$ has
$m\geq2$ preimages, use the same coefficient. Its right Kan
value is the homotopy limit of the star consisting of $m$
arrows $0\to M$, namely $(\Omega M)^{m-1}$; its left Kan
value is again zero. Taking $M$ to be the suspension of the
coefficient unit rules out this case. Hence every $r$ is
bijective. Conversely a finite discrete bifibration computes
both Kan extensions fibrewise, and the canonical finite
biproduct equivalence identifies them naturally under the
fibre-transport bijections and every strict base change.
\end{proof}

No invertibility of a prime is implicit in I.8. Its fibres have
no isotropy groups. If one returns to the larger smooth class,
the distinction matters: for $BC_p\to*$, trivial
$\mathbb F_p$ coefficients have homotopy orbits with nonzero
homology in every nonnegative degree and homotopy fixed points
with nonzero cohomology in every nonnegative degree. They are
not equivalent. The standard periodic resolution has
differentials alternating $1-g$ and $1+g+\cdots+g^{p-1}$;
both become zero on the trivial module. If $p$ is invertible
in a field, averaging splits invariants and the norm gives
ambidexterity for a finite $p$-group. These group projections
are not in the maximum diagonal-stable class: their diagonals
fail smoothness whenever the group is nontrivial. Thus this
diagnostic does not enlarge or qualify I.8.


\subsection{A specified enhanced second stage}
\begin{claim}{I.9}{Proved}{The second stages for spectra and mod-$p$ coefficients}
Fix a Grothendieck universe $\mathcal U$, and use a larger universe for categories of $\mathcal U$-small objects. Let $\mathcal C$ be either $\operatorname{Sp}_{\mathcal U}$ or $D(\mathbb F_p)$, with its indicated presentable infinity-category enhancement. We number the stages by
\[
 W_{\mathcal C}^{(0)}=W_{\mathbb D_{\mathcal C}},\qquad
 W_{\mathcal C}^{(1)}=W_{\mathbb D_{\mathcal C}}^{\mathrm{lc}},
 \qquad
 \mathcal C_1=\mathcal S_{W_{\mathcal C}^{(1)}},\qquad
 W_{\mathcal C}^{(2)}=W_{\mathbb D_{\mathcal C_1}}^{\mathrm{lc}}.
 \tag{I.9.1}
\]
Here $\mathcal S_W$ is the accessible reflective localization of spaces at exactly the class $W$, and
$\mathbb D_{\mathcal E}(A)=\operatorname{Ho}\operatorname{Fun}(NA^{\mathrm{op}},\mathcal E)$.
The localizations in (I.9.1) exist and preserve finite products. Their constant-coefficient classes satisfy
$W_{\mathbb D_{\mathcal C_1}}=W_{\mathcal C}^{(1)}$.
For every map $u:X\to Y$ of small spaces,
\[
 u\in W_{\mathcal C}^{(2)}
 \quad\Longleftrightarrow\quad
 \pi_0(u)\text{ is bijective and }
 \Omega_xu\in W_{\mathcal C}^{(1)}.
 \tag{I.9.2}
\]
The loop condition is tested at one basepoint in every component of $X$.
Consequently
\[
 W_{\operatorname{Sp}}^{(1)}
   =W_{\operatorname{Sp}}^{(2)}=W_\infty,
 \tag{I.9.3}
\]
and
\[
 u\in W_{D(\mathbb F_p)}^{(2)}
 \quad\Longleftrightarrow\quad
 \begin{cases}
 \pi_0(u)\text{ is bijective},\\
 \pi_i(u)\text{ is an integral isomorphism for }i=1,2,\\
 H_*(u\langle2\rangle;\mathbb F_p)\text{ is an isomorphism}.
 \end{cases}
 \tag{I.9.4}
\]
For a pointed connected component, $X\langle2\rangle$ means the homotopy fibre at its chosen basepoint of $X\to\tau_{\leq2}X$. The map $u\langle2\rangle:X\langle2\rangle\to Y\langle2\rangle$ is induced by $u$ and the basepoints. The last two conditions of (I.9.4) are tested on every corresponding pointed component. There are no nilpotence or finite-type hypotheses. The second class is again accessible and product stable, so $\mathcal C_2=\mathcal S_{W_{\mathcal C}^{(2)}}$ realizes it as an ordinary constant-coefficient class. This construction can be iterated at every finite stage for these two starting coefficients. It is a specified realization by localized spaces; functorial construction on all bare derivators and all exact morphisms remains the separate Open assertion of I.7.
\end{claim}
\begin{proof}
We prove the action comparison first under an explicit localization hypothesis, and then construct the localizations required here.

\paragraph{The constant class of a product-preserving localization.}
Let $W\subseteq W_0$ be the exact class inverted by an accessible reflector
$L:\mathcal S\to\mathcal E\subseteq\mathcal S$, and suppose that $W$ is stable under product with every small space. Local objects are closed under limits. The unit maps give
\[
 X\times Y\longrightarrow LX\times Y\longrightarrow LX\times LY
\]
in $W$, and $LX\times LY$ is local. Thus the specified comparison
$L(X\times Y)\to LX\times LY$ is an equivalence. The point is local, so $L$ preserves finite products. This is the meaning of Cartesian localization used here.

For any local space $b$ and any space $a$, the mapping space $\operatorname{Map}(a,b)$ is local: the test against $w\in W$ is the test of $w\times a$ against $b$. In particular every mapping space between two objects of $\mathcal E$ is itself a local space. For constants with values $a,b\in\mathcal E$, the comparison attached to $u:X\to Y$ is the map on components of
\[
 \operatorname{Map}\bigl(Y,\operatorname{Map}_{\mathcal E}(a,b)\bigr)
 \longrightarrow
 \operatorname{Map}\bigl(X,\operatorname{Map}_{\mathcal E}(a,b)\bigr).
\]
It is bijective when $u\in W$. Conversely, take $a=*$ and let $b$ range over all local spaces. Bijectivity gives
\[
 \operatorname{Hom}_{\operatorname{Ho}\mathcal E}(LY,b)
 \xrightarrow{\ \sim\ }
 \operatorname{Hom}_{\operatorname{Ho}\mathcal E}(LX,b).
\]
Yoneda in the homotopy category makes the specified $Lu$ an isomorphism, hence an equivalence. Therefore $W_{\mathbb D_{\mathcal E}}=W$.
The category $\mathcal E$ is presentable, so its functor categories have pointwise limits and colimits and both Kan extensions. These give Der~1--4 for the displayed $\mathbb D_{\mathcal E}$. Local constancy means that every arrow of the shape is sent to an equivalence, so the corresponding full infinity-category is $\operatorname{Fun}(BA,\mathcal E)$. Reversing paths identifies the presheaf convention with the space-valued convention used in the calculation.

\paragraph{The specified action monad and both implications.}
For a connected pointed space $(X,x)$, evaluation
\[
 U_x:\operatorname{Fun}(X,\mathcal E)\longrightarrow\mathcal E
\]
is conservative, because every point of $X$ is connected to $x$ by an invertible path. It preserves all colimits, which are computed pointwise. Its left adjoint $F_x=x_!$ sends $a$ to the copower of $a$ by the path spaces out of $x$. Consequently its monad is
\[
 T_X(a)=U_xF_x(a)
       =L(\Omega_xX\times a)
       \simeq L(\Omega_xX)\times a.
 \tag{I.9.5}
\]
The multiplication is induced by concatenation of based loops, and the unit by the constant loop. The infinity-categorical Barr--Beck theorem applies: the left adjoint exists, conservativity was checked, and preservation of geometric realizations follows from pointwise colimits \citep[Theorem~4.7.3.5]{GHIJLurie2017}. Hence $\operatorname{Fun}(X,\mathcal E)$ is the category of algebras for this particular monad.

Let $u:(X,x)\to(Y,y)$ be based, with both spaces connected, and let $R=u^*$. The identity $U_xR=U_y$ and composition of left Kan extensions give a canonical equivalence $u_!F_x\simeq F_y$. The unit of $u_!\dashv R$, evaluated on a free object, is therefore
\[
 F_x(a)\longrightarrow R F_y(a).
\]
Applying $U_x$ gives exactly the monad map
\[
 L(\Omega_xX\times a)\xrightarrow{L(\Omega_xu\times1_a)}
 L(\Omega_yY\times a).
 \tag{I.9.6}
\]
These are the specified maps, not abstract equivalences between their endpoints. If $R$ is an equivalence, this adjunction unit is an equivalence for every $a$. Taking $a=*$ proves that $L(\Omega_xu)$ is an equivalence, hence that $\Omega_xu\in W$. Conversely, if $\Omega_xu\in W$, product stability makes every map (I.9.6) an equivalence. It is a map of monads because $u$ preserves loop concatenation and the constant loop. Equivalence of these monads induces an equivalence of their algebra categories, and the natural Barr--Beck comparisons identify that restriction with $u^*$. This proves sufficiency. In particular an unspecified Morita equivalence is not used to infer equivalence of loop objects.

The same argument applies to the homotopy-category definition of $W^{\mathrm{lc}}$. Restriction has a left Kan-extension adjoint. If the induced functor on homotopy categories is an equivalence, the unit and counit of that adjunction are isomorphisms there and therefore equivalences in the infinity-categories. The original restriction is consequently an equivalence. The converse is immediate.

Both $\varnothing$ and $*$ are local, since $W\subseteq W_0$. If $u$ misses a target component, the map of local systems that is $\varnothing\to *$ on that component and the identity elsewhere becomes an equivalence after restriction, so restriction is not conservative. If two source components map into one target component, a source local system taking values $\varnothing$ and $*$ on those two components cannot be a restriction. Thus an equivalence of restriction functors forces a component bijection. Once components correspond, apply the connected calculation to each. This proves (I.9.2), including empty spaces and arbitrary component sets.

\paragraph{The first-stage coefficient tests.}
The preceding evaluation argument also works for any presentable coefficient category before localization. For spectra its monad is
\[
 a\longmapsto\Sigma^\infty_+\Omega_xX\wedge a,
\]
and for $D(k)$ it is
$ a\longmapsto C_*(\Omega_xX;k)\otimes_k^{\mathbf L}a$.
The same free-object unit, tested respectively on the sphere and on $k$, proves that a specified restriction is an equivalence exactly when its specified loop-spectrum or loop-chain map is an equivalence. The zero object and a nonzero coefficient detect components as above.

For spectra, smashing a loop-spectrum equivalence with $H\mathbb Z$ gives an integral homology equivalence of loop spaces. Its degree-zero map is the free abelian group map on the specified map of component groups
$\pi_1X\to\pi_1Y$. A map between free abelian groups induced by a function of bases is invertible only when that function is bijective. Thus the fundamental-group map is an isomorphism. Translation by loops identifies all components; restriction to the identity components gives an integral homology equivalence
\[
 \Omega\widetilde X\longrightarrow\Omega\widetilde Y.
\]
The bar construction carries this to an integral homology equivalence of $\widetilde X\to\widetilde Y$. Indeed normalized chains on a bar construction have a tensor-length filtration with the suspension of the augmentation ideal in positive degrees; the associated tensor-power comparisons are quasi-isomorphisms and the filtration is finite in each total degree. All chain complexes here are bounded below and degreewise free over $\mathbb Z$. The simply connected universal covers therefore have a weak equivalence by homological Whitehead. Together with the component and fundamental-group isomorphisms this proves $u\in W_\infty$. The reverse implication follows by homotopy invariance. We have proved $W_{\operatorname{Sp}}^{(1)}=W_\infty$; its localized-space coefficient is all of $\mathcal S$. Formula (I.9.2) now gives $W_{\operatorname{Sp}}^{(2)}=W_\infty$.

For $k=\mathbb F_p$, the same degree-zero argument gives the integral fundamental-group isomorphism, since
$H_0(\Omega X;k)=k[\pi_1X]$ and the algebra map is the one induced by $u$. Translation reduces the remaining loop-homology condition to $\Omega\widetilde u$. For maps of simply connected spaces $v$, one has
\[
 H_*(v;k)\text{ is invertible}
 \quad\Longleftrightarrow\quad
 H_*(\Omega v;k)\text{ is invertible}.
 \tag{I.9.7}
\]
For precision, use normalized chains on a $1$-reduced simplicial model. In the cobar construction the desuspended reduced generators have positive degree, so the tensor-length comparison converges in each degree; over a field tensor powers preserve quasi-isomorphisms. The cobar-to-loop-chain comparison identifies its conclusion with the forward direction of (I.9.7). The reverse direction is the bar argument just used. This is the classical Adams comparison; a primary modern statement and proof are \citet[Theorem~1]{I9Rivera2022}, and the $1$-reduced comparison is also made explicit in \citet[Proposition~4 and proof of Theorem~41]{I9RWZ2022}. The degreewise bounded filtration, rather than a finite-type assumption, is what is needed. This recovers D.5's first-stage criterion with all maps retained.

\paragraph{Existence and exactness of the chosen localizations.}
For either coefficient category, the class detected by restriction of local systems is closed under all small colimits of arrows of spaces: $\operatorname{Fun}(-,\mathcal C)$ sends colimits of spaces to limits of infinity-categories, and a limit of equivalences is an equivalence. It also has two-out-of-three and contains equivalences. Thus it is strongly saturated in the sense relevant to localization. It is product stable by the canonical identity
\[
 \operatorname{Fun}(X\times Z,\mathcal C)
   \simeq\operatorname{Fun}\bigl(Z,\operatorname{Fun}(X,\mathcal C)\bigr).
\]
For spectra this class is already $W_\infty$, so the localization is the identity. We supply set-generation for the mod-$p$ class, rather than infer accessibility merely from its closure properties.

Choose an uncountable regular $\kappa\in\mathcal U$. Represent a map of spaces by a simplicial-set map $f:K\to L$ and use the functorial Kan replacement $Q=\operatorname{Ex}^{\infty}$. This replacement preserves subcomplexes and filtered unions, and each simplex in $QK$ depends on finitely many simplices in $K$: at a finite $\operatorname{Ex}$-stage it is a map from a finite subdivided simplex. For a vertex $a$, use the based-loop model
\[
 \Omega_aQK=\{(a,a)\}\times_{QK\times QK}(QK)^{\Delta^1}.
\]
It too preserves filtered unions of pointed subcomplexes and has finite support in the original complex. A complex with fewer than $\kappa$ simplices gives fewer than $\kappa$ loop simplices and loop chains over $\mathbb F_p$.

Suppose $f\in W_{D(\mathbb F_p)}^{(1)}$, and start with a submap $K_0\to L_0$ of size less than $\kappa$. For each old target vertex choose an ambient source vertex and a finite path connecting its image to that target vertex. Whenever two old source vertices have connected images in the old target, choose an ambient path connecting those source vertices. For each old source vertex $a$ and every cycle $z$ in the target loop chains, choose an ambient source cycle $w$ and an ambient chain $b$ satisfying $f_*w-z=db$. For every old source cycle $w$ and old target chain $b$ satisfying $f_*w=db$, choose an ambient source chain $c$ with $dc=w$. The component bijection and loop-chain quasi-isomorphism of $f$ provide all these witnesses. There are fewer than $\kappa$ tests, and every witness has finite support in $K$ and $L$. Adjoin those supports, their faces and degeneracies, and their images under $f$.

Repeat countably many times. Regular uncountability keeps the union $K_\infty\to L_\infty$ smaller than $\kappa$. Every component relation, cycle, or bounding chain in this union occurs at a finite stage, so the next stage supplies its required witness. The union is therefore again in $W_{D(\mathbb F_p)}^{(1)}$. Given fewer than $\kappa$ such small weak submaps of $f$, apply the same procedure to their union to obtain an upper bound. Thus all these small weak submaps form a $\kappa$-filtered system whose union is $f$. Filtered unions of simplicial sets compute the corresponding colimits of spaces; this follows equally from the finite-support sphere and homotopy tests in $Q$.

There is a $\mathcal U$-small set of representatives of the simplicial-set maps of size less than $\kappa$ in this class. The preceding construction expresses every weak map as a filtered colimit of these representatives, up to isomorphism. Since the class is strongly saturated and closed under such colimits, it is exactly the strongly saturated class generated by that set. \citet[Proposition~5.5.4.15]{GHIJLurie2009} now gives an accessible reflective localization whose inverted maps are exactly $W_{D(\mathbb F_p)}^{(1)}$. The product-stability proof and the first paragraph make it Cartesian. This completes the existence hypotheses used in (I.9.2).

\paragraph{Evaluation of the second mod-$p$ stage and iteration.}
Apply the first-stage criterion to $\Omega_xu$. Its component condition is exactly that $\pi_1(u)$ be an isomorphism, and its fundamental-group condition on the identity component is exactly that $\pi_2(u)$ be an isomorphism. Translation carries that condition to every loop component. The universal cover of the identity component of $\Omega_xX$ is canonically equivalent, as a space over $\Omega_xX$, to
\[
 \Omega(X\langle2\rangle).
\]
Indeed this map has simply connected source, lands in the identity component, and is an isomorphism on every homotopy group in degrees at least two; the fibre sequence of the Postnikov truncation gives exactly the universal-cover map. This construction is natural in the pointed map $u$. Apply (I.9.7) to the $2$-connected covers to convert the loop-cover homology condition into the last line of (I.9.4). This proves the stated formula.

To realize this second class, the same bounded-submap proof uses the finite-support conditions on components and based loops together with the quasi-isomorphism test on $C_*(\Omega^2_xQK;\mathbb F_p)$. Component translations show that these are exactly (I.9.2)'s conditions. The additional fundamental-group witnesses are finite loop and homotopy data, so the same regular cardinal and countable closure procedure apply. The second class is strongly saturated and product stable because it is itself a locally constant coefficient class for $\mathcal C_1$. Consequently its exact accessible Cartesian localization exists and realizes it as a constant class by the first paragraph.

More generally at each finite stage the mod-$p$ criterion uses component bijectivity, the finitely many lower homotopy-group isomorphisms, and the quasi-isomorphism of chains on the corresponding iterated based-loop space. All representatives and witnesses still have finite support. Induction therefore constructs the subsequent localized-space coefficient categories; spectra remain at $\mathcal S$. The universal property makes these choices unique up to equivalence for each specified class. It does not furnish an assignment on unspecified bare derivators or a functor on all their exact morphisms. In (I.9.1), the phrase ``second stage'' thus refers precisely to the second locally constant step, after the first class has been realized by the chosen localized-space derivator.
\end{proof}
